
\documentclass[10pt,a4paper]{article}
\usepackage[utf8]{inputenc}
\usepackage{graphicx}
\usepackage{geometry}
\usepackage{hyperref}
\usepackage{array}
\usepackage{longtable}
\usepackage{pgfgantt}

\geometry{margin=1in}

\title{\textbf{Evaluating Multimodal Alignment in ChatGPT and Gemini: An Empirical Study of Inter- and Intra-Model Consistency in Image Description and Generation}}
\author{Saverio Napolitano \\ Mid Sweden University \\ \texttt{sana2504@student.miun.se}}
\date{}

\begin{document}
\maketitle

\section*{Background, Motivation, and Purpose}
AI models have demonstrated powerful capabilities in connecting text and images. 
However, it remains unclear how well their \textit{descriptive} and \textit{generative} capabilities align both across different models (inter-AI model) and within a single model (intra-AI model).
Misalignments can undermine trustworthiness, reproducibility, and usability in creative, scientific, and industrial contexts.
Alignment with human judgment remains imperfect as well. 
This project aims to systematically analyze inter-AI model and intra-AI model alignment in captioning and generation tasks, combining human and automatic evaluation.

\section*{Research Questions}
The study focuses on two main research questions:
\begin{enumerate}
    \item How consistent, with respect to human preferences, are multimodal AI models in aligning descriptions and generations, both within the same model and across different models?
    \item How well do automatic multimodal evaluation metrics align with human judgments?
\end{enumerate}

\section*{Relation to Related Work}

Surveys on multimodal LLMs \cite{wu2023survey,yu2025survey} outline architectures and alignment strategies, while CLIP-based evaluation methods and extensions like CLIP-AGIQA \cite{tang2024clipagiqa} allow automated quality assessment and have become standard tools. Recent approaches propose LLM-based evaluators \cite{liu2023geval}.

This work builds on representation alignment theory (e.g., CLIP embeddings \cite{radford2021clip}) and cycle consistency principles, where mutual agreement between description and generation provides a measure of alignment.

Unlike prior work focusing solely on evaluation metrics or model development, this project emphasizes comparative analysis of multiple AI models and their intra-model consistency, combining statistical analysis of human preference data with automated multimodal metrics.


\section*{Research Method}

A dataset of 10--15 diverse images will be built as input. Images will be collected explicitly for this study, ensuring AI models have not been pre-trained on them. 
Images will include objects (food, games, buildings, road signs, flags), text, nature (rivers, trees, grass, sky), people. 

Prompt engineering will be used to elicit captions and generations from selected AI models. 
Possible prompts include: making the AI model aware of the task (captioning or generation), specifying the target audience (same AI model, different AI model or humans), and requesting explanations for choices made in captions or generations to gain insight into the model's reasoning process; revealing the purpose of the task (e.g., for a human to rank the reconstructed image, for another AI model to generate the image, for the same AI model to generate the image) to see how context influences outputs.

The selected AI models are ChatGPT and Gemini, chosen for their advanced multimodal capabilities and accessibility. 
They will be used to generate captions for the images and to create images from the generated captions. 
Each AI model will use both its own captions and those generated by the other AI model to produce images, allowing for a comprehensive analysis of inter-AI model and intra-AI model alignment.

The generated images will be evaluated using both automated metrics (CLIP-AGIQA) and human judgment. 
The reference image will be provided as well. Human evaluators will rank the quality and relevance of the generated images based on the reference image.
Evaluators will be recruited from university students of different fields, ensuring a diverse range of backgrounds and perspectives to enhance the reliability of the human judgment data.

Kendall's W will be used to assess inter-rater agreement among human evaluators, while Kendall’s Tau and Spearman correlation will measure the correlation between human rankings and automated metrics.
Statistical analysis will be performed to identify significant differences in alignment across models and within models.

\section*{Expected Results}
We expect that
\begin{enumerate}
     \item Inter-AI model alignment will show significant divergence due to architectural heterogeneity, while intra-AI model descriptive vs. generative outputs will diverge less but still systematically. 
     \item CLIP-based automatic metrics will partially but not fully capture human preferences, requiring complementary human evaluation.
\end{enumerate}

\section*{Time Plan}
\begin{center}
\begin{ganttchart}[
        time slot format=isodate,      % daily input
        x unit=0.11cm,                 % << adjust this to change width
        hgrid,
        vgrid,
        y unit title=0.5cm,
        y unit chart=0.6cm,
        title label font=\scriptsize\bfseries,
        title height=1,
        bar label font=\scriptsize,
        group label font=\scriptsize
    ]{2025-10-13}{2026-01-11}

    % 13 ISO weeks from 2025-10-13 to 2026-01-11:
    % 42,43,...,52,1,2
    \gantttitlelist{42,43,44,45,46,47,48,49,50,51,52,1,2}{7} \\

    \ganttbar{Proposal Review}{2025-10-13}{2025-10-19} \\
    \ganttbar{Dataset Preparation}{2025-10-20}{2025-11-02} \\
    \ganttbar{Model Experiments}{2025-11-03}{2025-11-16} \\
    \ganttbar{Evaluation (Auto + Human)}{2025-11-16}{2025-11-30} \\
    \ganttbar{Results Analysis}{2025-12-01}{2025-12-07} \\
    \ganttbar{Article Writing}{2025-11-03}{2025-12-13} \\
    \ganttbar{Peer Review}{2025-12-14}{2025-12-21} \\
    \ganttbar{Project Presentation}{2025-12-22}{2026-01-11}

    \end{ganttchart}
\end{center}

\bibliographystyle{unsrt}
\bibliography{references}

\end{document}
